\begin{table}[htbp]
\centering
\caption[Hochschild 理论的低次数解释]{Hochschild 理论的低次数解释。这里 \(k\) 为域，\(A\) 为非零结合含幺 \(k\) 代数，\(M\) 为 \(A\) 双模。}
\label{b4:tab:hochschild-low-degrees}
\begin{tabularx}{\textwidth}{@{}p{0.22\textwidth}>{\raggedright\arraybackslash}X@{}}
\toprule
对象 & 代数意义 \\
\midrule
\(\operatorname{HH}_0(A,M)\)
& \(M/[A,M]\)，消去左右作用之差 \\
\(\operatorname{HH}^0(A,M)\)
& \(\{\,m\in M\mid am=ma\;\text{对所有}\;a\in A\,\}\) \\
\(\operatorname{HH}_1(A)\)
& 当 \(A\) 交换时，等于 Kähler 微分模 \(\Omega_{A/k}\) \\
\(\operatorname{HH}^1(A,M)\)
& \(\operatorname{Der}_k(A,M)/\operatorname{Inn}_k(A,M)\) \\
\(\operatorname{HH}^2(A)\)
& 保持单位及固定还原识别的一阶乘法形变等价类 \\
\(\operatorname{HH}^2(A,M)\)
& 固定核 \(M\)、商 \(A\) 及双模作用的平方零代数扩张等价类 \\
\(\operatorname{HH}^3(A)\)
& 给定有限阶形变的延拓障碍类所在的空间；并非每个三次类都必须出现为障碍 \\
\bottomrule
\end{tabularx}
\end{table}
